% TOP_PROBLEM: {"problemId": "problem.the-lindeloef-hypothesis-on-the-growth-of-the-riemann-zeta-function", "canonicalTitle": "The Lindeloef Hypothesis on the Growth of the Riemann Zeta Function", "releaseRank": 81, "primaryDomain": "number_theory_arithmetic_geometry", "primaryDomainLabel": "Number theory & arithmetic geometry", "displayStatus": "open", "releaseStatus": "open"}
% !TeX program = lualatex
% ATTEMPT_STATUS: unresolved
% Research continuation by GPT_6_Astra_Ultra, 27 September 2026.
\documentclass[11pt]{article}
\usepackage[margin=25mm]{geometry}
\usepackage{fontspec}
\setmainfont{Latin Modern Roman}
\usepackage{amsmath,amssymb,mathtools}
\usepackage{unicode-math}
\setmathfont{Latin Modern Math}
\usepackage{xurl,hyperref}
\hypersetup{colorlinks=true,urlcolor=blue}
\setcounter{secnumdepth}{0}
\setlength{\parindent}{0pt}
\setlength{\parskip}{5pt}
\setlength{\emergencystretch}{3em}
\begin{document}
\section*{\texorpdfstring{The Lindeloef Hypothesis on the Growth of the Riemann Zeta Function}{The Lindeloef Hypothesis on the Growth of the Riemann Zeta Function}}\label{81}
\subsection{Definitions and mathematical statement}
Let $\zeta$ be the meromorphically continued Riemann zeta function. The Lindelöf assertion is
\[
 \forall\varepsilon>0\ \exists C_\varepsilon,T_\varepsilon>0\quad
 |\zeta(\tfrac12+it)|\le C_\varepsilon|t|^\varepsilon
 \quad(|t|\ge T_\varepsilon).
\]
Thus growth on the critical line is smaller than every fixed positive power, with a constant allowed to depend on that power. The statement is a pointwise growth bound, not merely a mean-value estimate.
\par\smallskip\textit{Formulation and concept references:} \href{https://londmathsoc.onlinelibrary.wiley.com/doi/10.1112/jlms.70376}{[S1]}.

\subsection{Short English statement}
Does the zeta function on the critical line grow more slowly than every fixed positive power of height?
\subsection{Research attempt}
\paragraph{Scope and checked context.}
This unresolved attempt by GPT-6 Astra Ultra was prepared on 27 September
2026. Its main result is a self-contained reduction from bounds for every
fixed even moment to the pointwise statement, with the dependence of the
quantifiers made explicit. It then examines the off-diagonal terms which
prevent a short Dirichlet-polynomial calculation from proving those bounds.
None of these reductions is claimed to be new.

Florea's survey [E1], an accessible author version of [S1], distinguishes
unconditional bounds for smaller moments from results for larger moments
assuming the Riemann hypothesis. Matiyasevich's abstract [S3] explicitly
leaves convergence and exchanges of summation unjustified in its proposed
argument. The withdrawal at [S5] concerns version 48 of Unterberger's work;
the landing page [S4] now points to a later version. Neither an old
withdrawal nor the existence of a replacement establishes correctness of
the later claim. No claimed proof in these sources is used below.

\paragraph{An elementary derivative bound to control the width of a peak.}
Put $Z(t)=\zeta(1/2+it)$. For an integer $M\geq1$, summation by parts first
in $\operatorname{Re}s>1$ gives
\[
 \zeta(s)=\sum_{n\leq M}n^{-s}+\frac{M^{1-s}}{s-1}
       -s\int_M^\infty\{u\}u^{-s-1}\,du.                 \tag{81.1}
\]
For example the tail is $s\int_M^\infty(\lfloor u\rfloor-M)u^{-s-1}du$;
substituting $\lfloor u\rfloor=u-\{u\}$ proves the identity. The integral
in (81.1) converges locally uniformly for $\operatorname{Re}s>0$, and so
gives continuation there except at $s=1$. Differentiation under the integral
is justified on compact subsets by an integrable bound involving
$(1+\log u)u^{-1-\sigma}$ with $\sigma>0$.

Fix $T\geq3$, take $M=\lceil T\rceil$, and let $T/2\leq t\leq3T$.
Bounding the finite sum and the integral at $s=1/2+it$ yields
\[
 |Z(t)|\ll T^{1/2},\qquad |Z'(t)|\ll T^{1/2}\log T\ll T. \tag{81.2}
\]
For the derivative, the sum contributes at most
$\sum_{n\leq M}(\log n)/\sqrt n\ll\sqrt M\log(2M)$.
The differentiated integral contributes at most a constant times
$M^{-1/2}+T M^{-1/2}\log(2M)$. The middle term of (81.1) is smaller,
because $|s-1|\gg T$. These deliberately weak bounds are sufficient for
the reduction. In particular they do not assume Lindelöf or a sharp
subconvexity estimate.

\paragraph{Lemma: all fixed even moments imply the pointwise hypothesis.}
Assume the following family of estimates:
\[
 \forall k\in\mathbb Z_{>0}\ \forall\delta>0,\qquad
 \int_1^X|Z(t)|^{2k}\,dt\ll_{k,\delta}X^{1+\delta}.
                                                               \tag{81.3}
\]
Then Lindelöf holds.

\emph{Proof.} Let $t_0\in[T,2T]$ and $A=|Z(t_0)|$. Choose an absolute
$C$ such that the derivative bound in (81.2) is $CT$ on $[T/2,3T]$.
If $A=0$ there is nothing to prove. Otherwise on
$|t-t_0|\leq h=A/(2CT)$ the reverse triangle inequality gives
$|Z(t)|\geq A/2$. By increasing $C$, (81.2) ensures this interval lies
in $[T/2,3T]$. Thus
\[
 \int_{T/2}^{3T}|Z(t)|^{2k}dt
       \geq \frac{A^{2k+1}}{2^{2k}CT}.
\]
The assumed moment bound gives
\[
 |Z(t_0)|\ll_{k,\delta}T^{(2+\delta)/(2k+1)}.             \tag{81.4}
\]
Given the desired exponent $\varepsilon>0$, set $\delta=1$ and choose
one fixed integer $k$ so large that $3/(2k+1)<\varepsilon$.
The constants are then fixed before $T$ tends to infinity. Dyadic intervals
cover positive heights, and $\zeta(\overline s)=\overline{\zeta(s)}$
handles negative heights. This proves the stated pointwise bound.
\hfill$\square$

Conversely Lindelöf implies (81.3): use its pointwise exponent
$\delta/(2k)$, raise to the power $2k$, and integrate, absorbing the
bounded initial segment. Hence this moment family and the catalogue
statement are equivalent. The proof does not let $k$ grow with $T$ inside
an estimate whose implied constant is uncontrolled. It chooses $k$ only
after fixing the target exponent.

\paragraph{Why finitely many moments do not complete this route.}
With only one moment (81.4) retains a positive exponent. This is not
merely an aesthetic loss. Fix an integer $K\geq1$ and choose
$a=1/(K+1)$. At widely separated heights $T_j=4^j$, construct triangular
bumps of height $T_j^a$, width comparable to $T_j^{a-1}$, and slopes of
magnitude $T_j$. Set a function $g$ equal to these bumps and zero elsewhere.
The supports are disjoint for large $j$, and
\[
 \int_{\text{one bump}}|g(t)|^{2k}dt
       \asymp_k T_j^{(2k+1)a-1}.
\]
For every $k\leq K$ that exponent is at most $K/(K+1)<1$.
Summing the geometric sequence shows $\int_1^X|g|^{2k}\ll_K X$
for those finitely many moments, while $g(T_j)=T_j^a$ violates a
subpower bound. The example also satisfies a derivative bound of the
size used in (81.2), apart from harmless corners which can be smoothed.
It is not an analytic zeta-like function and is not a counterexample to
Lindelöf. It shows exactly why these finitely many integral estimates
and this continuity information do not force the required conclusion.

Similarly, Markov's inequality applied to one moment only gives
\[
 \operatorname{meas}\{t\in[T,2T]:|Z(t)|>T^\eta\}
       \ll_{k,\delta}T^{1+\delta-2k\eta}.                \tag{81.5}
\]
A small exceptional measure is compatible with peaks of positive height.
To eliminate every peak one needs information about its minimum width,
as in the lemma, and moments of arbitrarily high fixed order.

\paragraph{An exact polynomial calculation and its missing estimate.}
For a finite sum $D_N(t)=\sum_{n\leq N}a_n n^{-it}$, write
$D_N(t)^k=\sum_{m\leq N^k}b_m m^{-it}$, where $b_m$ is the sum of the
products $a_{n_1}\cdots a_{n_k}$ over $n_1\cdots n_k=m$.
Finite expansion and integration give the exact identity
\[
 \int_T^{2T}|D_N(t)|^{2k}dt
 =T\sum_m|b_m|^2+
 \sum_{m\ne n}b_m\overline{b_n}
       \int_T^{2T}e^{-it\log(m/n)}dt.                   \tag{81.6}
\]
Every off-diagonal integral has magnitude at most
$\min(T,2/|\log(m/n)|)$. For neighboring large indices $m,n$ of size
$N^k$, the logarithmic spacing is of size $N^{-k}$. When $N^k\gg T$,
oscillation alone need not make their individual integrals small. Discarding
them on grounds that $m\ne n$ would be false: exact orthogonality on a
finite interval is unavailable for these logarithmic frequencies.

For $a_n=n^{-1/2}$, the coefficients are
$b_m=d_{k,N}(m)/\sqrt m$, where $d_{k,N}$ counts restricted ordered
factorizations. The diagonal is then positive and accessible. Controlling
the sum over unequal products is a separate arithmetic problem, involving
close products with many factorizations. The familiar polynomial length
near $T^{1/2}$ in approximate-functional-equation approaches becomes
length near $T^{k/2}$ after taking the $k$th power. For large fixed $k$
this highlights the difficulty rather than proving a moment estimate.
No approximate functional equation or error estimate is being silently
used to identify the finite polynomial with zeta here; (81.6) is the
exact test calculation on its own.

\paragraph{Verification and first unresolved gap.}
The local script checks restricted-product coefficient counts and the
diagonal identity for small $N,k$ using exact rational arithmetic. It also
checks the moment-to-peak exponents and the bump exponents for several
fixed moment cutoffs. Those checks validate algebraic bookkeeping, not
the analytic hypothesis (81.3). Formula (81.1), its differentiated bounds,
and the peak-width implication are justified in the text.

The outcome is a quantitative proof that the complete moment family would
settle the pointwise conjecture, together with two explicit reasons why
weaker information does not: a finite list of moments permits narrow
peaks, and unequal products in (81.6) do not disappear. The first missing
input is an unconditional estimate (81.3) for every fixed $k$ (or another
argument controlling every peak directly). Neither source titles, a
formal summation exchange, nor a mean-value heuristic provides it.
The Lindelöf hypothesis remains \textbf{UNRESOLVED}.

\subsection{Sources}
\begin{itemize}
\item[S1]Alexandra Florea, A survey of moment bounds for zeta(s): From Heath-Brown's work to the present, Journal of the London Mathematical Society (2026).
\url{https://londmathsoc.onlinelibrary.wiley.com/doi/10.1112/jlms.70376}.
\textit{Formulation source.}
\item[S2]Analogues of the Lindelöf Hypothesis for the Barnes multiple zeta function and related problems.
\url{https://arxiv.org/abs/2606.02407}.
\textit{Further reference; not a proof endorsement.}
\item[S3]An approach to the Lindelöf Hypothesis for Dirichlet L-functions.
\url{https://arxiv.org/abs/2602.05731}.
\textit{Further reference; not a proof endorsement.}
\item[S4]Pseudodifferential arithmetic, Riemann and Lindelöf hypotheses.
\url{https://arxiv.org/abs/2208.12937}.
\textit{Further reference; not a proof endorsement.}
\item[S5]Pseudodifferential arithmetic, Riemann and Lindelöf hypotheses (withdrawn v48).
\url{https://arxiv.org/abs/2208.12937v48}.
\textit{Further reference; not a proof endorsement.}

\item[E1]Alexandra Florea, \emph{A survey of moment bounds for
$\zeta(s)$: from Heath-Brown's work to the present}, arXiv:2509.20335.
\url{https://arxiv.org/abs/2509.20335}.
\url{https://arxiv.org/html/2509.20335}.
\textit{Author version of [S1]. Inspected for the distinction between
unconditional and RH-conditional moment results. The all-moments implication
used here is proved explicitly in the attempt. Consulted 27 September 2026.}
\end{itemize}
\end{document}
